\chapter{Introduction}
\label{ch:intro}

MolDStruct SF/MD simulates what happens to a biomolecule, or any small system of light atoms, when it is hit by an
intense femtosecond \emph{optical} laser pulse.
The field of the pulse tears electrons out of the atoms by tunnelling, the sample charges up and explodes, and the
ions fly apart in directions that still carry information on the original structure.
The model follows the charge state of every atom, the electrons that are released, and the motion of the nuclei in one
simulation.

It is part of the MolDStruct toolbox and is built on GROMACS~4.5.4~\cite{pronk2013}, extended with an ionization
module.
Everything GROMACS does still works the same way; the ionization model is switched on with \code{mdrun -ionize} and
controlled by \opt{sfmd-*} options in the \file{.mdp} file.
Its sibling, MolDStruct MC/MD (\code{gromacs-mc}), treats X-ray pulses with a Monte Carlo of the inner-shell
processes; the two share the modified force field and, since November 2026, the trapped-electron cloud.

\section{The model in one page}

\begin{description}
  \item[Atoms] are classical particles with an integer charge.
    The charge is that of the ion's valence shell being stripped (H up to 1+, C 4+, N 5+, O 6+, S 6+), and it enters
    the force field.
  \item[The laser pulse] is a Gaussian in time with a carrier wave; its peak intensity follows from the pulse energy,
    duration and focal spot (Chapter~\ref{ch:adk}).
  \item[Tunnel ionization] is the Ammosov--Delone--Krainov (ADK) rate~\cite{adk1986,tong2002} of the instantaneous
    field for each atom's current charge state, turned into a probability per MD step and sampled (Chapter~\ref{ch:adk}).
  \item[The force field] is modified so that ionized atoms lose their chemical bonds: bonded terms fade out with the
    charge, bonds break irreversibly, and the Coulomb interactions between bonded neighbours are restored
    (Chapter~\ref{ch:ff}).
  \item[Released electrons] either escape or are trapped by the growing positive charge (optional,
    \opt{sfmd-electrons = pbgrid}).
    The trapped ones form a hot electron cloud, a continuous density in thermal equilibrium with the ions: a
    Poisson--Boltzmann problem solved on a three-dimensional grid every 0.2\,fs (Chapter~\ref{ch:pb}).
    The cloud screens the ions, pulls on them, decides which new electrons escape, and evaporates as the sample expands.
  \item[The laser acts on the cloud] by lowering its escape barrier and by heating it through inverse bremsstrahlung
    (Chapter~\ref{ch:laser}).
  \item[Electron--ion collisions] in the cloud ionize further, and after the pulse can recombine
    (Chapter~\ref{ch:coll}).
\end{description}

\section{When the cloud matters}

Without the cloud (\opt{sfmd-electrons = none}, the default) every electron leaves the instant it is freed, and the MD
sees bare ions.
For lysozyme (1960 atoms) at 850\,nm and 15\,fs that is a fair description at $10^{17}$\,W/cm$^2$, where the cloud is
gone before the pulse peak; at $10^{15}$\,W/cm$^2$ about half the electrons stay trapped, the ions reach only 0.3 of
their cloud-free kinetic energy, and the fast ions fly 6.5$^\circ$ away from their cloud-free directions
(Appendix~\ref{app:validation}).

\devnote{The cloud and everything built on it (the laser barrier and heating, impact ionization, recombination, the
temperature floor, long range, the direction autostop, the schedule) live on the \code{pb-electrons} branch of sf-md.
They were tested on lysozyme at $10^{15}$--$10^{17}$\,W/cm$^2$, 850\,nm, 15\,fs, but not yet against experiment or a
reference code; the results tables in Appendix~\ref{app:validation} are the evidence there is.}

\section{This manual}

Chapter~\ref{ch:start} gets a run going.
Chapters~\ref{ch:adk}--\ref{ch:coll} describe the physics, Chapter~\ref{ch:output} the output, and
Chapters~\ref{ch:mdp} and~\ref{ch:env} list every option and environment variable.
The appendices cover the numerics of the cloud, parallel runs, the tests behind the model, and version and compatibility
notes.
The code is in \file{src/kernel/sfionize.c} (ADK, the pulse, the run control) and \file{src/kernel/sfelectrons.c} (the
cloud), with the force field in \file{src/gmxlib/bondfree.c} and \file{nonbonded/nonbonded.c}.
